\section*{結言}
今回の8件では、文書編集や行政の答弁作成支援のように、生成AIが既存業務の工程へ入り込む動きが目立った。モデルの能力を競う発表だけでなく、誰がどの環境で使えるか、作業時間や品質にどのような効果があるかが、導入判断の論点になっている。

同時に、モデル行動仕様の改定、オープン音声モデルの普及、推論コストや資金調達を巡る議論も続く。ただし、今回の収集はGrok自身が\textit{Partial}と評価しており、一次投稿が観測窓外にある話題や未確認の二次情報が多い。とりわけ財務数値、Grok Botの成長率、ロボット試作報道は、Xで共有された主張と確定事実を区別して読む必要がある。
